% !TeX spellcheck = ru_RU
\chapter{Аналитическая часть}

В данной части приводятся основы устройства и работы протокола HTTP и подход к обслуживанию множества клиентов сервером отдачи статического содержимого с использованием мультиплексирования ввода-вывода посредством системного вызова \texttt{select()} и предварительно созданного набора процессов-обработчиков (worker'ов).

\section{Протокол HTTP}

HTTP (HyperText Transfer Protocol)~---~протокол клиент-серверного взаимодействия прикладного уровня для получения ресурсов, используемый в распределенных гипертекстовых системах~\cite{rfc9112}. В данном протоколе соединение и запросы к серверу инициируются получателем, сервер формирует ответ и отправляет его получателю~\cite{mdn_http_overview}.

\subsection{Структура HTTP-запроса}

HTTP-запрос состоит из трех основных частей:
\begin{enumerate}
\item стартовая строка (request-line)~---~метод, целевой ресурс, версия протокола разделенные одиночными пробелами;
\item заголовки (headers)~---~ноль или более строк в формате \texttt{field-name : field-value};
\item тело сообщения (body)~---~необязательная часть, которая обычно отсутствует в \texttt{GET} или {HEAD} запросах.
\end{enumerate}


\subsection{Методы HTTP-запросов}

Всего определено 8 методов HTTP-запросов~\cite{rfc9110}:
\begin{itemize}
	\item \texttt{GET}~---~получение содержимое целевого ресурса;
	\item \texttt{HEAD}~---~то же что и \texttt{GET}, но без передачи сервером тела ответа;
	\item \texttt{POST}~---~выполнение специфичной для данного целевого ресурса обработки содержимого запроса;
	\item \texttt{PUT}~---~замена всех текущих представлений целевого ресурса содержимым запроса;
	\item \texttt{DELETE}~---~удаление всех текущих представлений целевого ресурса;
	\item \texttt{CONNECT}~---~установление соединения с сервером, определяемым целевым ресурсом;
	\item \texttt{OPTIONS}~---~получение разрешенных методов и заголовков для взаимодействия с целевым ресурсом;
	\item \texttt{TRACE}~---~получение запроса в том виде, в котором он был доставлен серверу.
\end{itemize}

Обязательными для реализации серверами общего назначения являются только методы \texttt{GET} и \texttt{HEAD}~\cite{rfc9110}.

\subsection{Структура HTTP-ответа}
HTTP-ответ также состоит из трех основных частей, аналогичных частям HTTP-запроса:
\begin{enumerate}
	\item строка состояния (status-line)~---~версия протокола, код состояния, необязательное текстовое описание кода состояния разделенные одиночными пробелами;
	\item заголовки (headers)~---~ноль или более строк в том же формате, что и заголовки HTTP-запроса;
	\item тело сообщения (body)~---~необязательная часть, которая отсутствует в \texttt{HEAD} запросах и может также отсутствовать в {GET} запросах.
\end{enumerate}

\subsection{Коды состояния НТТР}
Выделяют 5 классов кодов состояния HTTP~\cite{rfc9110}:
\begin{itemize}
	\item \texttt{1xx}~---~информационный, указывает на промежуточный ответ, для сообщения о состоянии соединения или ходе выполнения запроса до отправки окончательного ответа;
	\item \texttt{2xx}~---~успешный, указывает на то, что запрос был успешно получен, понят и обработан;
	\item \texttt{3xx}~---~перенаправление, указывает на то, какие дальнейшие действия необходимо предпринять клиенту для выполнения запроса;
	\item \texttt{4xx}~---~ошибка клиента, указывает на то, что запрос клиента выглядит ошибочным;
	\item \texttt{5xx}~---~ошибка сервера, указывает на то, что сервер допустил ошибку в ходе выполнения данного запроса, или не способен обработать данный метод.
\end{itemize}

В описываемом в данной работе сервере используются коды состояния: \texttt{200 OK}, \texttt{403 Forbidden}, \texttt{404 Not Found}, \texttt{405 Method Not Allowed}.

\section{Мультиплексирование ввода-вывода}

Системный вызов \texttt{select()} позволяет программе отслеживать несколько файловых дескрипторов, ожидая, пока один или несколько из них не станут готовыми для какой-либо операции ввода-вывода. Файловый дескриптор считается готовым, если возможно выполнить соответствующую операцию ввода-вывода без блокировки. Таким образом системный вызов \texttt{select()} позволяет осуществлять мультиплексирование ввода-вывода~\cite{linux_select}.

Основным недостатком системного вызова \texttt{select()} является то, что он позволяет отслеживать только файловые дескрипторы, номера которых меньше \texttt{FD\_SETSIZE} (1024).

\section{Использование пула процессов-обработчиков}

Архитектура с использование многопроцессной обработки запросов предполагает, что родительский процесс создает набор дочерних процессов, посредством системного вызова \texttt{fork()}~\cite{linux_fork}, до получения каких-либо запросов, после чего процессы-потомки обрабатывают получаемые запросы, а процесс-предок отслеживает их состояние и создает новый дочерний процесс, в случае, если один из дочерних процессов завершился.

Преимуществом данного подхода является изолированность процессов-обработчиков друг от друга. К недостаткам можно отнести: высокое потребление памяти, так как каждый обработчик представляет собой отдельный процесс; сложность балансировки нагрузки; проблема thundering-herd, когда несколько обработчиков просыпаются для выполнения системного вызова \texttt{accept()}~\cite{linux_accept}.


\section*{Вывод}

В данном разделе были рассмотрены основы устройства и работы протокола HTTP и подход к обработке множества клиентов с помощью мультиплексирования ввода-вывода системным вызовом \texttt{select()} и предварительно созданного набора процессов-обработчиков.

\clearpage